\section{Lower Bounds for Monotone Circuits}
\label{sec:monotone-degree-LB}

Recall that $\calM_n(\ell)$ denotes all Boolean monotone $\ell$-slice
functions on $n$ bits: all Boolean functions
$f: \{0,1\}^n \rightarrow \{0,1\}$ that output $0$ on all inputs of
Hamming weight less than $\ell$ and $1$ on all inputs of Hamming
weight larger than $\ell$. There is no restriction on the output for
inputs of Hamming weight $\ell$ and we have
$|\calM_n(\ell)| = 2^{\binom{n}{\ell}}$.
%
Further, recall that $\calM_n(\ell, s, t) \subseteq \calM_n(\ell)$ is
the class of monotone Boolean $\ell$-slice functions
$f: \set{0,1}^n \rightarrow \set{0,1}$ for which there is a
(not necessarily monotone) Boolean circuit
$C_f^\mon: \set{0,1}^n \rightarrow \set{0,1, \bot}$ of size $s$ such
that
\begin{enumerate}
\item for all $\ell$-slice inputs $\alpha \in \binom{[n]}{\ell}$ it
  holds that if $C^\mon_f(\alpha) \neq \bot$, then
  $C^\mon_f(\alpha) = f(\alpha)$, and
\item $C^\mon_f(\alpha) = \bot$ on at most $t$ inputs
  $\alpha \in \binom{[n]}{\ell}$.
\end{enumerate}

It is very convenient to work with slice functions as we have a handle
on their monotone circuit complexity: by \cref{lem:slice} their
monotone circuit size is the same as their ordinary circuit size up to
a polynomial size increase. Hence we do not need to worry whether the
functions needed for the reduction have small monotone circuits, as
long as we are working on a slice only. 

The proof of the monotone lower bound is an adaption of the argument
used to prove
\cref{lem:restriction}. % We prove the monotone lower bounds by
% adapting the proof of \cref{lem:restriction}.
The idea is to work over
the $\ell$th slice and disregard all other inputs.  By
\cref{lem:slice} we can implement our selector circuits by small
monotone circuits. We then also need to take care of the negations in
the $\oplus$-circuit. We push the negations down until they either hit
a gate in $Y$ or a selector circuit. We create a set $\overline Y$
gates, which we can think of as the negation of the gates in $Y$ and
also create negated selector circuits (on the $\ell$th slice). By
doing so we can now get rid of the last negations by appropriately
connecting the appropriate circuits. The following corollary of
\cref{lem:slice} will be useful to us.

\begin{claim}
  \label{clm:to-slice}
  Let $C$ be a Boolean circuit on $n$ input bits of size $s$. Then,
  for $\ell \in [n]$, there is a monotone Boolean circuit $C^\mon$ of
  size $2s + \poly(n)$ computing the $\ell$-slice function that is
  equal to $C$ on the $\ell$-slice.
\end{claim}

\begin{proof}
  Let $\calT_{\ge \ell}$ be the threshold function that outputs $1$ if
  and only if the Hamming weight of an input $\alpha \in \set{0,1}^n$
  is at least $\ell$. Connect the output of $C$ by an \emph{and} gate
  to a circuit computing $\calT_{\ge \ell}$. The output of this
  circuit is then connected by an \emph{or} gate to the output of a
  circuit computing $\calT_{> \ell}$. Let us denote this new circuit
  by $C'$.

  The circuit $C'$ clearly outputs $1$ whenever the input is of
  Hamming weight larger than $\ell$. Furthermore, on the $\ell$-slice
  it is equal to $C$ because $\calT_{\ge \ell}$ outputs $1$ while
  $\calT_{> \ell}$ outputs $0$. Finally the output is $0$ if the
  Hamming weight is less than $\ell$ because the output of both
  threshold functions is $0$.

  Clearly the size of the circuits computing the threshold functions
  is $\poly(n)$. We apply \cref{lem:slice} to conclude that there is a
  monotone circuit $C^\mon$ computing the same function as $C'$ of
  size $2s + \poly(n)$.
\end{proof}

Before stating the following lemma we need to adapt some terminology
to the monotone setting. Observe that $\circuit_s^\mon(f)$ is a
restriction of $\circuit_s(f)$. Let $\tau$ be such that
$\restrict{\circuit_s(f)}{\tau} = \circuit^\mon_s(f)$. This allows us
to naturally extend $k$-determined restrictions to the monotone
setting: a restriction $\rho$ is a $k$-determined restriction for
$\circuit^\mon_s(f)$ if the restriction $\rho \tau$ is a
$k$-determined restriction for $\circuit_s(f)$. Similarly we can
extend $m$-independence to the monotone setting. This will later allow
us to use \cref{lem:circuit-to-poly} even though we are working with
the monotone formula.


\begin{lemma}
  \label{lem:restriction-monotone}
  For all $k, \ell, m, n, t \in \N$ satisfying $m \le 2^n$, and any
  explicit bipartite graph $G = (U, V, E)$ such that $|U| = 2^n$,
  $|V| = m$ and all $u \in U$ are of degree $\deg(u) \le k$, the
  following holds.  There is a constant $C > 0$, depending on the
  explicitness of $G$, such that for all
  $s \ge C \cdot m \cdot n^C \cdot k^C$ and any
  $f \in \calM_n(\ell, s/10, t)$ there is a natural $m$-independent
  $k$-determined variable substitution $\rho$ for the formula
  $\circuit^\mon_s(f)$ such that
  \begin{align*}
    \gout_{s, \alpha}(Y) =
    \begin{cases}
      1, &\text{if~} |\alpha| > \ell,\\
      0, &\text{if~} |\alpha| < \ell,\\
      f(\alpha), &\text{if~} |\alpha| = \ell \text{~and~} C_f^\mon(\alpha) \neq \bot,\\
      \oplus_{i \in N(\alpha)} y_i, &\text{otherwise},\\
    \end{cases}
  \end{align*}
  for $\gout_{s, \alpha}$ and $Y$ as in \cref{def:determined}.

  Furthermore, the formula $\restrict{\circuit^\mon_s(f)}{\rho}$ is
  over $O(t \cdot k + m)$ variables.
\end{lemma}

\begin{proof}
  This proof is an adaptation of the argument of the proof
  \cref{lem:restriction}. Let us describe the natural $m$-independent
  $k$-determined restriction $\rho$ for the formula
  $\circuit_s^\mon(f)$.
  
  As in the proof of \cref{lem:restriction} we have gates that act as
  Boolean variables. But instead of having a single set $Y$ of
  variables we now have two sets $Y$ and $\overline Y$, each of size
  $m$. We think of the variables in $\overline Y$ as the negations of
  the variables in $Y$ and ensure this by applying the appropriate
  variable substitution for all $\alpha \in \set{0,1}^n$ and $i \in [m]$.

  According to \cref{clm:to-slice} we may assume that the circuit
  $C_f^\mon$ computes a monotone $\ell$-slice function in both outputs
  $C_{f,1}^\mon, C_{f,2}^\mon$ for a mild increase in size;
  $|C_f^\mon| \le s/5 + \poly(n) \le s/4$ for $s$ large enough. Recall
  that the first output of $C_f^\mon$ indicates whether the second
  output bit is equal to $f$ on the $\ell$-slice. Let
  $\overline C^\mon_{f,1}$ be the negation of $C_{f,1}^\mon$ on the
  $\ell$-slice. In other words,
  $\overline C^\mon_{f,1}(\alpha) = \lnot C_{f,1}^\mon(\alpha)$ if
  $\alpha$ has Hamming weight $\ell$, and
  $\overline C^\mon_{f,1}(\alpha) = C_{f,1}^\mon(\alpha)$ otherwise.

  The monotone circuit $C_f^\mon$ is of size at most $s/4$ and hence
  according to \cref{lem:slice} there is a monotone circuit of size
  $s/2 + \poly(n) \le 5s/8$ computing
  $\overline C^\mon_{f,1}(\alpha)$.

  We restrict the formula such that a part of the circuit is
  equivalent to $C_f^\mon$ and another part is equal to
  $\overline C^\mon_{f,1}$. Note that the size of these two circuits
  is at most $7s/8$ by above discussion.

  Recall that because $G(\set{0,1}^n, Y, E)$ is explicit, there are
  circuits $\sel_1,$ $\sel_2, \ldots, \sel_m$, each of size $\poly(n,k)$, where
  each $\sel_i$ computes, given an input $\alpha \in \set{0,1}^n$,
  whether the vertex $y_i \in Y$ is a neighbor of the vertex $\alpha$.
  Let $\overline \sel_i = \lnot \sel_i$ and denote by $\sel_i^\mon$
  (respectively $\overline \sel_i^\mon$) the circuit obtained by
  applying \cref{clm:to-slice} to $\sel_i$ (to $\overline \sel_i$
  respectively). By the guarantees of \cref{clm:to-slice} all these
  $2m$ circuits are of size $\poly(n,k)$.

  We restrict the formula such that a part of the circuit computes the
  functions
  \begin{align}
    \sel_1^\mon, \ldots, \sel^\mon_m,
    \overline\sel^\mon_1, \ldots, \overline\sel^\mon_m \eqperiod
  \end{align}
  From these $\ell$-slice selector circuits we can then define
  selector circuits that take $C_f^\mon$ into account. Namely, we
  connect $\sel_i^\mon$ by an \emph{and} gate to the output of
  $\overline C_{f,1}^\mon$ to obtain the circuit $\sel_{i}^{\prime\mon}$ and
  similarly connect $\overline \sel^\mon_i$ by an \emph{or} gate to
  $C_{f,1}^\mon$ to obtain the circuit $\overline \sel_i^{\prime\mon}$. 

  Finally, we also put each variable $y_i$ and $\bar y_i$ onto the
  slice by the same construction used in the proof of
  \cref{clm:to-slice}: connect the variable $y_i$ (respectively
  $\bar y_i$) by an \emph{and} to the threshold circuit
  $\calT_{\ge \ell}$ and connect this circuit in turn by an \emph{or}
  gate to a $\calT_{> \ell}$ threshold circuit to obtain $y_i^\mon$
  (respectively $\bar y_i^\mon$).  It is well-known \cite{Valiant84,
    BW06, Goldreich2020} that threshold circuits have montone circuits
  of size $\poly(n)$ and we can thus restrict the formula such that a
  part of the circuit computes $y_i^\mon$ and $\bar y_i^\mon$.
  
  Finally we connect $y^\mon_i$ by an \emph{and} gate to the selector
  circuit $\sel_i^{\prime\mon}$. Note that this circuit is equal to an
  $\ell$-slice function. As we will see later this ensures that the
  whole circuit outputs an $\ell$-slice function.
  %
  We connect the circuits $\bar y^\mon_i$ similarly: connect
  $\bar y^\mon_i$ by an \emph{or} gate to the negated selector circuit
  $\overline \sel_i^{\prime\mon}$. Again, the output of this circuit
  is equal to an $\ell$-slice function.

  Equally inportant is that these circuits behave well on the
  $\ell$-slice. Indeed it can be checked that the positive circuit, on
  input $\alpha \in \set{0,1}^n$, outputs
  $\sel_i^{\prime\mon}(\alpha) \land y_i$ while the negative circuit
  outputs $\overline \sel_i^{\prime\mon}(\alpha) \lor \bar y_i$. On
  the $\ell$-slice these functions are the negation of eachother,
  which we are going to use in the following.
  
  We need to construct a monotone circuit for the \emph{xor} of
  $\sel_i^{\prime\mon}(\alpha) \wedge y_i$ for $i$ from $1$ to $m$, on
  $\ell$-slice inputs $\alpha$. We take a standard $O(m)$-size
  $\oplus$-circuit and monotonize it by pushing all negations in it
  down using De Morgan's law until they reach one of the
  $\oplus$-circuit's inputs $\sel_i^{\prime\mon} \wedge y_i$.
  Whenever the negation of $\sel_i^{\prime\mon}(\alpha) \wedge y_i$ is
  needed, we do one last step of De Morgan and replace it by
  $\overline \sel_i^{\prime\mon}(\alpha) \vee \bar{y}_i$.

  To ensure that the circuit outputs $f(\alpha)$ whenever
  $C^\mon_f(\alpha) \neq \bot$, we connect the two outputs of
  $C_f^\mon$ by an \emph{and} gate and connect this gate by an
  \emph{or} gate to the output of the \emph{xor} circuit. This
  completes the description of the restriction on the structure
  variables. A depiction of the resulting circuit can be found in
  \cref{fig:monotone}.
  \begin{figure}
    \centering
    \includegraphics[width=\linewidth,height=\textheight,keepaspectratio]{fig/circuit-restricted-monotone.pdf}
    \caption{A depiction of the monotone circuit, where
      $\widetilde\oplus$ is the $\oplus$ circuit with the negations
      pushed down.}
    \label{fig:monotone}
  \end{figure}
  We ensure that $s$ is large enough so that above circuit can be
  described by the formula.

  Note that the constructed circuit always outputs a monotone
  $\ell$-slice function: as the monotonized $\oplus$-circuit is
  non-constant, we see that if all inputs to the circuit are $0$, it
  outputs $0$ and if all inputs are $1$, it outputs $1$. This, in
  particular, implies that the circuit outputs $0$ (respectively $1$)
  if the input is below (respectively, above) the $\ell$-slice and
  hence the entire circuit computes a monotone $\ell$-slice function.

  It can be easily checked that the described restriction is
  $m$-independent and $k$-determined. In order to prove the
  furthermore part, we need to reduce the number of evaluation
  variables. This can be achieved analogous to the proof of
  \cref{lem:restriction} and we thus omit it here.
\end{proof}

Let us prove our degree lower bound for monotone circuits, restated
here for convenience.

\MonotoneDegree*

\begin{proof}[Proof of \cref{thm:main-monotone-degree}]
  As in the proof of \cref{thm:main-degree}, we use the graphs from
  \cref{thm:expander}, with $U = \{0,1\}^n$,
  $k = O_{\gamma}\big((n \log r)^{1 + 1/\gamma}\big)$, and
  $|V| \le k^2 r^{1+\gamma}$ for parameters $\gamma > 0$ and
  $r \le 2^n$. We apply \cref{lem:restriction-monotone} with above
  graph and $t=2^n$ along with $C_f^\mon = \bot$ to obtain, for
  $s \ge m \cdot \poly(n, k)$, an appropriate natural $m$-independent
  $k$-determined variable substitution $\rho$ for
  $\circuit^\mon_s(f)$. In particular $\rho$ satisfies
  \begin{align*}
    \gout_{s, \alpha}(Y) =
    \begin{cases}
      1, &\text{if~} |\alpha| > \ell,\\
      0, &\text{if~} |\alpha| < \ell,\\
      \oplus_{i \in N(\alpha)} y_i, &\text{otherwise},\\
    \end{cases}
  \end{align*}
  for $\gout_{s, \alpha}$ and $Y$ as in definition
  \cref{def:determined}.

  Recall that there is a restriction $\tau$ such that
  $\circuit_s^\mon(f) = \restrict{\circuit_s(f)}{\tau}$ and we can
  thus apply \cref{lem:circuit-to-poly} with $\tau\rho$ to conclude
  that if there is an SoS refutation of
  $\restrict{\circuit^\mon_s(f)}{\rho}$ in degree $d$, then there is a
  degree $d \cdot k$ SoS refutation of the system of polynomial
  equations computing
  \begin{align}
    \set{
    \bigoplus_{i \in N(\alpha)} y_i = f(\alpha)
    \mid
    \alpha \in \binom{[n]}{\ell}
    }
    \eqperiod
  \end{align}
  %
  As the graph $G$ is a strong expander, we can apply
  \cref{thm:sos-xor} to get an SoS degree lower bound of $\Omega(r)$
  for above system of equations. By above connection this gives an
  $\Omega(r/k)$ degree lower bound for the
  $\restrict{\circuit^\mon_s(f)}{\rho}$ formula and hence also for the
  unrestricted formula.

  Regarding the parameters, as in the proof of \cref{thm:main-degree}
  we choose $m = s/\poly_\gamma(n)$. Repeating the
  calculations from the aforementioned proof we obtain that
  $r \ge s^{1/(1+\gamma)}/\poly_\gamma(n)$. Thus by choosing
  $\gamma$ small enough such that $\frac{1}{1 + \gamma} > 1 - \eps/2$
  and $s$ large enough such that the final
  $\poly_\gamma(n) \le s^{\eps/2}$ we obtain the claimed degree lower
  bound of $\Omega_\eps(s^{1-\eps})$.
\end{proof}

As in the non-monotone case, we can also obtain size lower bounds for
functions that almost have a circuit of size $s$.

\MonotoneSize*

\begin{proof}
  Analogous to the proof of \cref{thm:main-size}.
\end{proof}